\documentclass[a4paper]{article}
\usepackage{amsfonts}
\usepackage{amsmath}
\usepackage{amssymb}
\usepackage{graphicx}  
\usepackage{cancel}

\newcommand{\degree}{^{\circ}}
\newcommand{\cis}{\operatorname{cis}}
\newcommand{\Bgsin}{\operatorname{Bgsin}}
\newcommand{\Bgcos}{\operatorname{Bgcos}}
\newcommand{\Bgtan}{\operatorname{Bgtan}}
\newcommand{\limit}{\lim\limits}

\begin{document}
  
\noindent \large Auteur: Rune Suy \\
\noindent \large Studierichting: Informatica\\
\noindent \large Oefeningenreeks 3C nr. 8\\

\medskip

\normalsize

Zij $a, b \in \mathbb{R}$, $a<b$. Gebruik de middelwaardestelling om te bewijzen dat: \\

$\exists c \in ]a, b[ : (n+1) c^n = \sum_{i=0}^{n} a^i \cdot b^{n-i}$.\\

\bigskip

Gegeven: $a, b \in \mathbb{R}$, $a<b$\\

Te bewijzen: $\exists c \in ]a, b[ : (n+1) c^n = \sum_{i=0}^{n} a^i \cdot b^{n-i}$\\

Bewijs:\\

Stel $f(x) = x^{n+1}$\\

$f$ is continu op $[a, b]$ en differentieerbaar op $]a, b[$\\

$\Rightarrow_{Lagrange} \exists c \in ]a,b[ : f'(c) = \dfrac{f(b) - f(a)}{b-a}$\\

$\Leftrightarrow (n+1)c^n = \dfrac{f(b) - f(a)}{b-a} = \dfrac{b^{n+1} - a^{n+1}}{b-a}$\\

Nog te bewijzen: $b^{n+1} - a^{n+1} = (b-a)\displaystyle\sum_{i=0}^{n} a^i \cdot b^{n-i}$\\

Bewijs: \\

$b^{n+1} - a^{n+1} = (b-a)\cdot(b^n + ab^{n-1} + ... + a^{n-1}b + a^n)$\\

$= (b-a)\sum_{i=0}^{n} a^i \cdot b^{n-i}$\\

$QED$


\begin{thebibliography}{999}
%\bibitem{a} Referentie
\end{thebibliography}
\end{document} 
